\documentclass[11pt]{article}
\usepackage[letterpaper,margin=1in]{geometry}
\usepackage[T1]{fontenc}
\usepackage{lmodern,microtype}
\usepackage{amsmath,amssymb,booktabs,array,longtable}
\PassOptionsToPackage{svgnames}{xcolor}
\usepackage{graphicx,xcolor,tikz}
\usepackage{caption}
\usepackage[numbers,sort&compress]{natbib}
\usepackage[colorlinks=true,allcolors=blue!45!black,breaklinks=true]{hyperref}
\hypersetup{pdftitle={The Challenge of Deriving Signal from Synthetic Legal Environments: Evidence from Harvey's Legal Agent Benchmark},pdfauthor={John J. Hughes, III}}
\urlstyle{same}
\captionsetup{font=small,labelfont=bf}
% Drafting markers. A line starting with % TODO is visible only in the source.
% Colored macros mark provisional prose. For a clean PDF, redefine \draft,
% \checkcite, and \checknum to return their argument unchanged.
\newcommand{\draft}[1]{\textcolor{red}{#1}}
\newcommand{\checkcite}[1]{\textcolor{blue}{#1}}
\newcommand{\checknum}[1]{\textcolor{orange}{#1}}
% A criterion quoted verbatim from Harvey's LAB task files: an indented block with a
% label above and the source below, so quoted text is never mistaken for the author's.
\newenvironment{labquote}[1]{\def\labquotesource{#1}\begin{quote}\small
{\footnotesize\textsc{Quoted from Harvey LAB}}\par\nobreak}
{\par\nobreak{\raggedleft\footnotesize --- \labquotesource\par}\end{quote}}
% Heading of one sampled criterion: number, Harvey's task title, criterion id.
\newcommand{\labitem}[3]{\subsubsection*{#1. #2 (#3)}}
% BEGIN GENERATED LAB REVIEW NUMBERS
% Counts from the human review of sampled Harvey LAB criteria.
\newcommand{\labTasks}{52}
\newcommand{\labCriteria}{2{,}858}
\newcommand{\labSolFlagged}{616}
\newcommand{\labOpusBlind}{650}
\newcommand{\labOpusFinal}{564}
\newcommand{\labFlagged}{835}
\newcommand{\labFlaggedPct}{29.2}
\newcommand{\labSample}{25}
\newcommand{\labSampleTasks}{19}
\newcommand{\labDefective}{16}
\newcommand{\labArguable}{3}
\newcommand{\labNotDefective}{6}
\newcommand{\labObjective}{9}
\newcommand{\labEnvItems}{14}
\newcommand{\labEnvTasks}{10}
\newcommand{\labAnyDefect}{20}
\newcommand{\labNoDefect}{4}
\newcommand{\labBothFlagged}{10}
\newcommand{\labBothDefective}{9}
\newcommand{\labSingleFlagged}{15}
\newcommand{\labSingleDefective}{7}
\newcommand{\labDefectiveLB}{383}
\newcommand{\labDefectivePoint}{534}
\newcommand{\labDefectiveShare}{13.4}
\newcommand{\labObjectiveLB}{171}
\newcommand{\labObjectivePoint}{301}
\newcommand{\labObjectiveShare}{5.9}
\newcommand{\labEnvLB}{319}
\newcommand{\labEnvShare}{11.1}
\newcommand{\labEnvPoint}{468}
\newcommand{\labAnyLB}{524}
\newcommand{\labAnyPoint}{668}
% END GENERATED LAB REVIEW NUMBERS
\setlength{\emergencystretch}{2em}
% Flush paragraphs. A gap marks a new paragraph; no line is indented.
\usepackage[skip=0.6em]{parskip}
\title{The Challenge of Deriving Signal from Synthetic Legal Environments:\\Evidence from Harvey's Legal Agent Benchmark}
\author{John J. Hughes, III\\Independent researcher\\\url{https://www.johnjhughesiii.com/research/}}
\date{Working draft --- October 7, 2026}
\begin{document}
\maketitle
\begin{abstract}
Reinforcement learning (RL) techniques, in domains that can offer objectively verifiable rewards, have led to incredible and noteworthy recent gains in the capabilities of frontier AI models. But not all domains are equally verifiable. Many high-level legal judgment tasks are difficult to verify, and model capability gains on those tasks have not been nearly as noteworthy as model gains in domains that are susceptible to automated verification like math and coding.

Existing benchmarks used to evaluate high-level legal tasks rely largely on human-written rubrics that test model output in response to fictitious, AI-generated environments. But an AI-assisted audit of one of the most widely used such benchmarks, checked by a practicing litigator's review of a random sample of the flagged criteria, found evidence of hundreds of erroneous rubrics in just one subpart of the benchmark, many of them in task environments that are themselves misspecified. A separate AI-assisted citation check found apparently fabricated case citations in the task environments, and a deep dive into one motion-to-dismiss task found planted issues and documents that do not resemble real lawyer work product. The errors in these environments illustrate the challenges facing even the top data providers and research labs trying to improve frontier model capabilities, at a time when many previously challenging benchmarks and tasks have become saturated.
\end{abstract}

\section{Introduction}

Most recent legal benchmarks have sought to test higher-level judgment tasks, which usually lack an objective ground truth. The leading vertical legal AI company, Harvey AI Corporation, released Legal Agent Benchmark (LAB), in May 2026. Now widely cited by frontier research labs as a key evaluation of agentic legal capabilities for frontier models, LAB uses AI-generated task environments to test models on pre-specified rubrics. Harvey has publicly disclosed that it works with Mercor and Snorkel, leading third-party data providers, to conduct expert review and verification of every generated LAB task \citep{harveyResearchLab}. Harvey open-sourced the benchmark (except for a held-out test set), and invited practicing lawyers to validate and improve it by reviewing existing tasks and auditing rubrics. \citep{harveyLAB} I took up that invitation and did a deep-dive review of one of the substantive litigation tasks (drafting an issues list for a motion to dismiss opposition brief) and an audit of all the litigation rubrics. This analysis underscored the challenges involved in creating suitable reinforcement learning environments that can advance the capabilities frontier in the legal domain---and illustrated why extracting signal from real-world litigation data could be more effective than existing approaches \citep{lfbpaper}.


\section{Audit of LAB's litigation rubrics}

Several LAB rubrics appeared to contain anomalous results, so I conducted an audit of all rubrics on litigation and dispute resolution tasks (my area of specialty for 18 years). I supplied GPT-6 Sol and Claude Opus 5.5 with examples of the kinds of rubrics that I consider defective and asked them to iterate through each task to audit the rubrics using the high-level standards I supplied (which allowed them to raise issues that were not specifically on my list).

Of the 52 tasks and 2,858 rubrics included in the litigation and dispute resolution category, GPT-6 Sol flagged 616 criteria as defective or arguably defective, Claude Opus 5.5 flagged 564 (with the ability to review GPT-6 Sol's findings). Between them, the models flagged 835 criteria, leaving 2,023 unflagged by either report. Of 345 overlapping flags, both reviewers called 83 problematic and 135 arguable; 127 received different levels of concern. Opus was allowed to see Sol's findings when reaching its own conclusions, so agreement in the final reports is not independent confirmation \citep{labaudit}.

To validate these findings and ascertain a statistically valid estimate of the number of rubrics that are defective, I reviewed a simple random sample of \labSample{} of the \labFlagged{} flagged criteria, drawn with a fixed seed from \labSampleTasks{} tasks (Appendix~\ref{app:labreview}). I treated a rubric as defective if, applied as written to the task as given, it would fail a response a careful litigator would consider correct or pass one that is wrong. \labDefective{} of the \labSample{} were defective, and I assessed \labObjective{} of those to be unambiguous objective errors where the accuracy of the rubric did not depend on my exercise of professional judgment but there was indisputably an error or problem in the task of some kind. Another \labArguable{} were arguable and \labNotDefective{} were not defective.

These results demonstrate that at least \labDefectiveLB{} of the \labFlagged{} flagged criteria are defective (one-sided 95\% bound; point estimate \labDefectivePoint{}), or more than \labDefectiveShare\% of all \labCriteria{} criteria in the practice area; at least \labObjectiveLB{} are unambiguous objective errors. Agreement between the two AI auditors was informative: \labBothDefective{} of the \labBothFlagged{} sampled criteria flagged by both were defective, against \labSingleDefective{} of the \labSingleFlagged{} flagged by only one.

Examples of defective rubrics, drawn from the sampled criteria in Appendix~\ref{app:labreview}, included instances where the rubric misstates the law (e.g., requiring the answer that is the opposite of settled law in the jurisdiction at issue), mathematical or numerical errors (e.g., the task environment contained mathematically impossible documents, or contained a rubric that required models to assume that one of two conflicting numbers was accurate while ignoring the conflicting number), or instances where I assessed the rubric's recommended approach to contravene widely respected professional norms.


I separately counted defects in the task environment relevant to the rubric (as opposed to defects that made the rubric's grading criteria themselves incorrect). \labEnvItems{} of the \labSample{} criteria related to defective task environments, for example, documents that did not sum to the stated figures; these span \labEnvTasks{} of the \labSampleTasks{} sampled tasks, and only \labNoDefect{} of the \labSample{} criteria had neither a defective nor an arguable criterion and came from an environment without a defect. This means that at least \labEnvLB{} flagged criteria (more than \labEnvShare\% of all criteria) come from such environments.

These errors are highly significant relative to the margins that separate frontier models on LAB's scoring criteria. On Artificial Analysis's Harvey LAB-AA leaderboard, Muse Spark 1.3 leads with a criterion pass rate of 95.5\%, followed by Kimi K3 at 94.6\%, Claude Fable 5.1 at 93.3\%, Claude Opus 5.5 at 91.2\%, and GPT-5.6 Luna at 87.9\% \citep{labAAboard}. If criteria elsewhere in LAB are defective at anything like the rate we found (more than \labDefectiveShare\% of the criteria in this practice area), the rate of defective rubrics is as large as or larger than the reported gaps between leading models on criterion pass rate, and it is unclear how much of those gaps the defects could explain. All-pass scoring compounds the problem, because a single defective criterion can fail an otherwise acceptable task.

The flaws in these environments raise questions about whether general-purpose data vendors have the qualifications or ability to prepare verified and accurate RL environments for legal tasks. Harvey retained leading data vendors to review and verify the tasks and rubrics \citep{harveyResearchLab}. But the verification services offered by these third-party vendors did not reliably detect defects like rubrics that misstated the law. These general-purpose data vendors provide a variety of different kinds of reinforcement learning environments, many of them for nonlegal tasks, and the companies utilized do not purport to be specialized in legal work. Some companies in this space hire freelance lawyers at contract attorney rates to verify tasks and rubrics. At an earlier phase of frontier model development, that approach may have provided useful signal. But as the capabilities of frontier models have improved, and benchmarks have become more ambitious, the results here suggest that this kind of freelancer data could have defects that end up undercutting the evaluation or training signal the dataset is intended to provide. Even sophisticated buyers doing frontier AI research in the legal domain may find it difficult to catch these errors without effectively reproducing much of the work that the data vendors are paid to do. This suggests that further improvement in the legal domain might require efforts to increase the quality of the data used to evaluate and train models on legal tasks.

The challenges involved in developing and maintaining acceptable benchmarks for legal tasks call to mind some of the challenges frontier labs have had developing new benchmarks for coding tasks. In February 2026, OpenAI publicly criticized SWE-bench Verified (which OpenAI had helped develop), noting that its audit revealed problems involving test requirements and task specification \citep{swecritique}. The rapid capability growth has meant that errors that once might have been unnoticeable now show up in the form of frontier models failing tasks incorrectly because the model is smarter than the benchmark. That dynamic could require greater investment in evaluations to keep pace with the strength of the underlying models.

\section{Cite Check of LAB's Task Environment and Rubrics}
\label{sec:citecheck}
Because the rubric audit turned up instances where the rubrics contained hallucinated case law (suggesting that the external experts who were responsible for verification may have relied on inaccurate AI-generated responses to prepare the answer key), I asked AI agents to conduct a separate citation check of all the litigation and dispute resolution tasks to assess the frequency of AI hallucinations in the task environments or rubrics. A citation check is a standard pre-filing procedure where lawyers check every authority cited in a brief, regulatory submission, or other important document to ensure that the authority is cited accurately, contains any quotation attributed to it, and supports the legal proposition for which it is cited.

A separate AI-assisted check of every case citation in the \labSampleTasks{} sampled environments against CourtListener, with spot checks, found problems in 45 of 102 citations to purported real authority: 32 real cases cited for holdings, quotations, or facts they do not contain, 7 cases that could not be found, 5 garbled citations to real cases, and 1 citation pointing to a different case \citep{labreview}.

After excluding errors in case law that were intentionally planted and then tested in rubrics (e.g., an error of law in an opponent's submission that the model was supposed to catch), the AI identified 436 errors that were apparently not intended to be part of the task environment \citep{labcitecheck}:

\begin{table}[ht]\centering\small
\begin{tabular}{lrr}\toprule
Finding & All citations & Case citations\\\midrule
Verified & 2{,}754 & 306\\
Verified, minor defect (pin cite, loose paraphrase, wrong year) & 346 & 54\\
Intentional error tested by the rubric (excluded) & 31 & 6\\\midrule
\multicolumn{3}{@{}l}{\textit{Unintentional errors}}\\
\quad Cited for a proposition it does not support & 246 & 105\\
\quad Not found (likely fabricated) & 76 & 60\\
\quad Garbled citation to a real authority & 58 & 10\\
\quad Quotation not found in the authority & 40 & 31\\
\quad Reporter citation belongs to a different case & 16 & 12\\
\quad Total unintentional errors & 436 & 218\\\midrule
Source text not accessible & 51 & 6\\\midrule
Total checked & 3{,}618 & 590\\
\bottomrule\end{tabular}
\caption{AI citation check of the \labTasks{} LAB litigation task environments. This table shows the AI findings.}
\label{tab:citecheck}
\end{table}

\draft{[[Placeholder: To add results of attorney audit of the results, the hallucinated cases issue was identified only on 10/4/2026, the day before our original intended publication date. We delayed publication to understand how widespread the issue was, and initial spot checks suggest there are a meaningful number of hallucinated citations in the task environments so we need to present that with some validation so the issue is properly quantified.]]}

\section{A motion-to-dismiss case study}
\label{app:lab}

In addition to auditing the accuracy of all 2,858 rubrics, I did a deep dive review of one substantive litigation task. This review highlighted a different but equally important issue: the use of synthetic or AI-generated data may result in task environments that are unrealistic in a variety of subtle ways that may not be obvious to nonexperts but that are striking and clear to practicing attorneys. The task \texttt{analyze-counterparty-motion-to-dismiss} asks AI agents to prepare a list of issues for the opposition brief to a motion to dismiss, on a fictitious software service dispute.

I observed several ways in which the synthetically generated task environment seemed unrealistic and could potentially train models on a distorted view of what genuine practice looks like. \citep{labtask}

\paragraph{Planted issues.} A number of the legal disputes in the task environment are obviously ``planted'' in a way that feels artificial and ultimately presents issues in a way that is meaningfully different from how real-world legal disputes manifest. For example, the contract the parties are suing over has a forum selection clause that requires any lawsuit to be filed in Travis County, Texas. The plaintiff literally quotes this entire provision in the complaint but nevertheless files suit in the Northern District of Georgia and offers no explanation for why they believe the forum selection clause they quoted does not apply or otherwise permits the suit in Georgia. It of course happens regularly that a lazy or incompetent litigant files suit in the wrong forum; a competent lawyer potentially could decide to ignore the clause and hope that their adversary overlooks the issue. But in both of those scenarios, the forum selection clause would not be explicitly flagged in the complaint as it is here; a real version of this legal task would require the lawyer to dig through a number of documents and realize that buried in one of them is a forum selection clause that allows the lawyer to have the case moved to Texas.

This is one example, but a significant number of the tasks in these environments were similarly planted in a way that's artificial and unlike how disputes around these issues arise in real cases. Environments like this can distort an evaluation of a model's true capabilities, since simply looking for the obviously planted issues will satisfy the benchmark without testing whether the model would be able to identify this issue from more realistic source material akin to what it will see in real cases.

\paragraph{Environment Documents Violate Court Rules.} Many of the documents in these environments do not resemble real lawyer work product. The environment provides counterparty litigation documents in docx format, and the native files from the git repository often cannot be opened, although agents and various CLI tools can recover the document text. In real litigation, lawyers rarely share their docx files with their opponent; the opponent's memorandum of law is ordinarily filed with the court as a PDF. The motion papers also lack a proper caption, include a stale table of contents, contain markdown formatting artifacts that render in Word as junk characters even in files that can be opened normally, lack citations for propositions that generally require citations, and are single spaced rather than double spaced in violation of the applicable court rules.

The motion that the agent is being asked to respond to is so far below professional norms that you would want any analysis of the motion to include consideration of whether the judge would strike the motion completely, or potentially discipline or sanction the counsel that filed that deficient work product. That is not included in the rubric for this task because the quality control issues appear to be unintentional, not part of the test for the agent. But it is a genuine gap: if a document like this were filed in real litigation, the quality of the work product could itself become a substantive issue in the litigation and thus would need to be flagged given the goal motivating the all-pass criterion, which was to ensure that the AI does not overlook any important issues \citep[section ``Evaluating a task'']{harveyLAB}.

There are valid reasons to prefer rubric-based reinforcement learning approaches. For example, while a benchmark like LegalForecastBench \citep{lfbpaper} (which asks models to predict how federal judges ruled on motions to dismiss) can test a model's legal reasoning and analytical ability over a narrow task, it does not test the model's ability to generate new work product, which is an important consideration for a company like Harvey that is developing legal agents to prepare and revise different kinds of legal work product. But reviewing these environments suggests that the current state of the art in synthetic data generation introduces drift on a variety of dimensions that could undermine the efficacy of the task environment as a training or evaluation signal. The model is in effect being tested on its ability to recognize a distinct environment, very unlike what it will encounter in production, and adhere to a stylized set of norms and rules that were artificially constructed in that environment but do not comport with the real norms or rules that govern real-world tasks. While creating LegalForecastBench illustrated that there are real challenges in working with real-world litigation data as well, for a variety of tasks that data has a lot of rich, implicit signal that data providers generating synthetic data are unable to capture.

\section{Conclusion}

This audit illustrated that AI research has a long way to go to bridge the gap between current training and evaluation environments used to train models, and the messy real-world tasks that models are asked to perform. LAB is notable as the first large-scale, open-source set of legal task environments that try to test high-level legal judgment tasks where there is often no objective ground truth. But LAB illustrates the challenges of creating effective environments, because even leading data vendors that took on the responsibility of expert verification of the LAB tasks were unable to create realistic task environments or accurate rubrics. Harvey has explicitly invited practicing attorneys to help collaborate on improving LAB, and building models that can automate more of the high-level legal work that lawyers do clearly will require careful collaboration to obtain training signal that can reliably improve the capabilities of, and distinguish between the performance of, increasingly capable frontier models.

\section*{Reproducibility, access, and AI use disclosure}
Code and audit materials are maintained in the \href{https://github.com/johnhughes3/LegalForecastBench}{LegalForecastBench repository}.

As I am a practicing attorney and do not have any formal training in artificial intelligence or machine learning, the use of AI models was essential to allow me to make a research contribution in this space that otherwise would have been beyond my capabilities. My personal contribution included auditing a sample of the LAB rubrics by hand to statistically validate the AI audit findings, conducting a deep dive into one of LAB's litigation tasks, and orchestrating AI coding agents that helped write and execute the code that ran the audit. I wrote most of this paper personally, but discussed the results at several points with AI agents and considered AI-generated analysis as an input into my thinking as I drafted. The major exception to this was the AI audit findings, the citation check in Section~\ref{sec:citecheck}, and other places where numerical results are included; I relied on AI agents to produce those findings, to write code that performs the necessary calculations, and to prepare the LaTeX source summarizing those results. I used several different models from Anthropic, OpenAI, SpaceXAI, and Google to check the work. I also met with experienced researchers in the space to validate the experimental design and results, and obtain feedback on the draft. I am grateful to several AI researchers in industry positions for helpful discussions. This version of the paper is intended as a working draft, and feedback from researchers, lawyers, and others on how to improve this audit or make our findings more useful is welcome.


\clearpage

\begin{thebibliography}{99}

\bibitem{labaudit} GPT-6 Sol audit workers, Claude Opus 5.5 audit workers, Coordinating agents; individual provenance not fully recorded. Harvey LAB litigation and dispute resolution audit: GPT-6 Sol and Claude Opus 5.5. \emph{Local research artifact, pinned September 29, 2026}, 2026. \url{https://github.com/johnhughes3/LegalForecastBench/tree/58371a0e13494e591710e2eae3da8e428eebc359/docs/harvey-lab-audit/litigation-dispute-resolution}.

\bibitem{labreview} John Hughes. Human review of AI-flagged Harvey LAB litigation criteria: protocol, sample, and worksheet. \emph{Research artifact}, 2026. \url{https://github.com/johnhughes3/LegalForecastBench/tree/main/docs/harvey-lab-audit/litigation-dispute-resolution/human-review}.

\bibitem{labcitecheck} Claude Sonnet 5.5 and Claude Opus 5.5 citation-check agents, coordinated by the author. Citation check of the Harvey LAB litigation and dispute resolution task environments. \emph{AI-generated research artifact, not attorney-reviewed}, 2026. \url{https://github.com/johnhughes3/LegalForecastBench/tree/main/docs/harvey-lab-audit/litigation-dispute-resolution/citation-check}.

\bibitem{labtask} Harvey AI task authors; upstream authorship not individually listed. Pinned LAB task.json for analyze-counterparty-motion-to-dismiss. \emph{Primary online source}, 2026. \url{https://github.com/harveyai/harvey-labs/blob/1dd81403b2fbb60596f7aea3fcecafad7bf73143/tasks/litigation-dispute-resolution/analyze-counterparty-motion-to-dismiss/task.json}.

\bibitem{harveyResearchLab} Gabe Pereyra. How Harvey Built a Research Lab on a Budget. \emph{Sequoia Capital, Own Your Intelligence}, August 11, 2026. Slide at 4:11: ``Scale with expert data providers.'' \url{https://www.youtube.com/watch?v=MGouk8W51v0\&t=251s}.

\bibitem{harveyLAB} Niko Grupen, Gabe Pereyra, Julio Pereyra. Introducing Harvey's Legal Agent Benchmark. \emph{Primary online source}, 2026. \url{https://www.harvey.ai/blog/introducing-harveys-legal-agent-benchmark}.

\bibitem{swecritique} OpenAI. Why SWE-bench Verified no longer measures frontier coding capabilities. \emph{Primary online source}, 2026. \url{https://openai.com/index/why-we-no-longer-evaluate-swe-bench-verified/}.

\bibitem{labAAboard} Artificial Analysis. Harvey LAB-AA evaluation results. \emph{Leaderboard, accessed October 5, 2026}, 2026. \url{https://artificialanalysis.ai/evaluations/harvey-lab-aa}.

\bibitem{lfbpaper} John J. Hughes, III. Law as a Verifiable Domain: Using Real-World Judicial Outcomes to Evaluate AI Legal Reasoning. \emph{Working paper}, 2026. \url{https://www.legalforecastbench.org/paper/}.

\end{thebibliography}

\clearpage

\appendix
\section{Human review of AI-flagged LAB criteria}
\label{app:labreview}
I asked two frontier models, GPT-6 Sol and Claude Opus 5.5, to audit the \labTasks{} tasks in LAB's litigation tasks (specifically Harvey commit \texttt{1dd81403}). Table~\ref{tab:labaudit} reports how many criteria each pass flagged as defective or arguably defective. The combined row is Opus's later review of rubrics already flagged by either pass. A criterion ``flagged by either model'' was flagged at least once; that is \labFlagged{} of the \labCriteria{} criteria (\labFlaggedPct\%) \citep{labaudit}.

\begin{table}[ht]\centering\small
\begin{tabular}{llr}\toprule
Pass & What the model saw & Criteria flagged\\\midrule
GPT-6 Sol & The tasks, without Opus's results & \labSolFlagged{}\\
Claude Opus 5.5, first pass & The tasks, without Sol's results & \labOpusBlind{}\\
Claude Opus 5.5, combined & The rubrics flagged by either pass & \labOpusFinal{}\\\midrule
Flagged by either model & Share of all \labCriteria{} criteria & \labFlagged{} (\labFlaggedPct\%)\\
\bottomrule\end{tabular}
\caption{AI audit of the \labTasks{} litigation tasks at Harvey commit \texttt{1dd81403}.}
\label{tab:labaudit}\end{table}

In order to validate the AI findings and determine how many of the AI-flagged rubrics were genuinely defective, I reviewed a simple random sample of \labSample{} of the flagged rubrics, drawn with a fixed seed from \labSampleTasks{} tasks. I committed the protocol for that review before it began \citep{labreview}. For each rubric, I reviewed the grading rubric itself, the findings from both AI auditors, how LAB's AI judges graded two real model runs (GPT-6 Luna and Opus 5.5) on that rubric, and the relevant documents from the task environment. In a few cases I also did research into the relevant governing law. I marked a rubric as defective when it could produce an inaccurate grade, meaning: either the rubric could fail a response that a thoughtful litigator would consider correct, or the rubric could pass a response that is actually incorrect on the issue the rubric is testing.

During the review, I decided to distinguish between defects in the rubrics themselves and defects in the task environment. The line between the two was based on whether the defect could change the grade a model is given on the task. For example, one rubric wanted the work product to contain a specific metric sourced from the task environment, but the AI auditors pointed out that the task environment actually contained two contradictory numbers for that particular metric. Accordingly, an AI agent that relied on one of the numbers (which was not the one the rubric was expecting) would be failed incorrectly because of the underlying defect in the task environment. I treated that as a defect in both the rubric and the task environment. By contrast, a different task environment contained a notice of an OSHA regulatory violation from a state-level regulatory agency in Ohio that does not actually have jurisdiction over the violations charged. The rubric wanted work product to simply flag the existence of the charge in the work product. I coded that as a defect in the task environment but not as a defect in the rubric (because an AI agent handed that task could, and likely still would, pass the rubric by simply mentioning the charged violation).

Of the \labSample{} criteria, I graded \labDefective{} rubrics defective in ways that affect grading, including \labObjective{} that I considered to be unambiguous objective errors in which the record, the law, the cited authority, or the arithmetic contradicts the rubric. \labArguable{} were arguable defects, and \labNotDefective{} were not defective. \labEnvItems{} of the \labSample{} come from task environments that contain defects (\labEnvTasks{} of the \labSampleTasks{} tasks). Only \labNoDefect{} had neither a defective nor an arguable criterion and came from an environment without a defect (Table~\ref{tab:labreview}).

Since the sample was a draw without replacement from the \labFlagged{} criteria that AI agents had flagged as defective, the one-sided 95\% lower confidence bounds are: at least \labDefectiveLB{} flagged criteria are defective in ways that affect grading (point estimate \labDefectivePoint{}; the bound is \labDefectiveShare\% of all criteria in the practice area); at least \labObjectiveLB{} are unambiguous objective errors (point estimate \labObjectivePoint{}); and at least \labEnvLB{} come from task environments that contain defects (point estimate \labEnvPoint{}). At least \labAnyLB{} involve a defective criterion, a defective environment, or both (point estimate \labAnyPoint{}). In the sample, \labBothDefective{} of the \labBothFlagged{} criteria flagged by both auditors were defective, against \labSingleDefective{} of the \labSingleFlagged{} flagged by only one.

I did not attempt to estimate defects in criteria that were not flagged by the AI agents; it is possible there could be some additional defects that the agents did not flag. These estimates count criteria/rubrics, not the individual mistakes that give rise to the defects. In some cases, the defect was specific to the particular rubric at issue, but in other cases, a single error in a task environment could make several criteria defective. After I prepared an initial set of verdicts, I asked AI models (GPT-6.1 Sol and Opus 5.5) to stress test my answers and to draft additional arguments both for and against each call and a summary of those they would change. I changed a handful of verdicts after reflecting on these points, and did a second pass through the full set of my verdicts to ensure that they applied the principles I adopted during the review consistently to each call.

The audit of rubrics revealed instances where the task environment, and sometimes grading rubrics, appeared to rely on AI-generated hallucinated case law citations. As a result, I had AI agents (Opus 5.5) conduct a separate AI-assisted check of every case citation in the \labSampleTasks{} sampled task environments against CourtListener. One agent checked each task. A synthesis agent re-verified a sample of the findings, and I independently re-checked seven of them. This process found problems in 45 of 102 citations to real authority, including 7 cases that could not be found and 32 real cases cited for propositions they do not contain \citep{labreview}.

% BEGIN GENERATED LAB REVIEW TABLE
\begin{table}[ht]\centering\small
\begin{tabular}{lrrr}\toprule
Criterion verdict & Environment defective & Environment correct & Total\\\midrule
Defective & 10 & 6 & 16\\
Arguable & 2 & 1 & 3\\
Not defective & 2 & 4 & 6\\
\midrule
Total & 14 & 11 & 25\\
\bottomrule\end{tabular}
\caption{My verdicts on the 25 sampled criteria, by whether the task environment itself contains a defect.}
\label{tab:labreview}\end{table}
% END GENERATED LAB REVIEW TABLE

\subsection*{The sampled criteria}
Each entry has four parts: (1) the criterion quoted verbatim from Harvey's task file (that text is from Harvey's data vendors, not me, and in many cases the text appears to be AI-generated), (2) my verdict on the criterion and task environment, (3) results from a sample run of LAB using GPT-6 Luna and Claude Opus 5.5, to assess how LAB's AI graders (Claude Sonnet 4.6 and GPT-5.5) graded outputs from Luna and Opus on the specific criteria at issue, (4) my explanation of the verdict. Harvey task files are cited at commit \texttt{1dd81403}.


\labitem{1}{Assess Reasonableness of Staffing Levels on Litigation Invoice}{C-049}
\begin{labquote}{Harvey LAB, criterion C-049, \href{https://github.com/harveyai/harvey-labs/blob/1dd81403b2fbb60596f7aea3fcecafad7bf73143/tasks/litigation-dispute-resolution/assess-reasonableness-of-staffing-levels-on-litigation-invoice/task.json}{task file at commit \texttt{1dd81403}}}
\textbf{Calculates budget overage percentage as approximately 17\%.} PASS if the memo calculates the percentage by which the attorney fee subtotal (\$193,372.50) exceeds the approved budget ceiling (\$165,000) as approximately 17\% (exact figure is 17.2\%). FAIL if no percentage overage calculation is provided or if the percentage stated is materially incorrect (off by more than 2 percentage points).
\end{labquote}

\noindent\begin{tabular}{@{}>{\raggedright\arraybackslash}p{.25\linewidth}>{\raggedright\arraybackslash}p{.71\linewidth}@{}}
\toprule
\multicolumn{2}{@{}l@{}}{\textbf{My assessment}}\\
Criterion & Defective --- unambiguous objective error, deficiency in task environment\\
Task environment & Defective: the invoice's 117 time entries add up to \$215,680.00, while its summary and subtotal show \$193,372.50.\\
AI auditors' reasoning & Correct\\
\midrule
\multicolumn{2}{@{}l@{}}{\textbf{LAB's own grading of two model runs} (Harvey's judges, not my assessment)}\\
GPT-6 Luna & Sonnet 4.6 judge: fail; GPT-5.5 judge: fail\\
Claude Opus 5.5 & Sonnet 4.6 judge: pass; GPT-5.5 judge: pass\\
\bottomrule
\end{tabular}

\noindent\textbf{My analysis.} The rubric is testing the AI on a specific percentage calculation from an invoice, but the document in the environment is internally inconsistent. The invoice’s summary page and its subtotal row show fees of \$193,372.50 (17.2\% over the \$165,000 budget), but the 117 detailed time entries add up to \$215,680.00, which puts the overage at \$50,680, or 30.7\% (13.5 points from the rubric’s required 17.2\% answer and far outside the 2-point tolerance).

\labitem{2}{Categorize Document Production Set by Relevance and Privilege — Privilege Log and Relevance Classification Report}{C-012}
\begin{labquote}{Harvey LAB, criterion C-012, \href{https://github.com/harveyai/harvey-labs/blob/1dd81403b2fbb60596f7aea3fcecafad7bf73143/tasks/litigation-dispute-resolution/categorize-document-production-set-by-relevance-and-privilege/task.json}{task file at commit \texttt{1dd81403}}}
\textbf{ISSUE\_003: DOC\_012 identified as Partially Privileged.} PASS if DOC\_012 (chandrasekaran-meeting-notes-aug28.docx, Chandrasekaran's handwritten meeting notes from August 28, 2023) is classified as Partially Privileged. FAIL if classified as fully Privileged or Not Privileged.
\end{labquote}

\noindent\begin{tabular}{@{}>{\raggedright\arraybackslash}p{.25\linewidth}>{\raggedright\arraybackslash}p{.71\linewidth}@{}}
\toprule
\multicolumn{2}{@{}l@{}}{\textbf{My assessment}}\\
Criterion & Defective --- low-quality professional judgment, rubric is incorrect\\
Task environment & Correct\\
AI auditors' reasoning & Correct\\
\midrule
\multicolumn{2}{@{}l@{}}{\textbf{LAB's own grading of two model runs} (Harvey's judges, not my assessment)}\\
GPT-6 Luna & Sonnet 4.6 judge: pass; GPT-5.5 judge: pass\\
Claude Opus 5.5 & Sonnet 4.6 judge: pass; GPT-5.5 judge: pass\\
\bottomrule
\end{tabular}

\noindent\textbf{My analysis.} The rubric is wrong because it fails fully Privileged, which is the better answer here (even if Partially Privileged is defensible).

\begin{itemize}
\item The documents are handwritten lawyer notes from a fictitious meeting discussing problems with the R\&D Department. The headings on individual topics seem intended to create the impression that there is a discussion of business topics (which would not be privileged) followed by a separate discussion of legal topics (privileged), so from a superficial glance, the rubric does seem to match the topic headings. However, the bullets under each heading make clear that the entire meeting is just dedicated to the topic of the company’s desire to find a way to fire the head of R\&D without getting sued. The business topics (which look to be nonprivileged) are discussions of how the R\&D department is a bit of a disaster, such that management wants to fire the head. The legal sections then involve discussion of the fact that the R\&D head has some sort of environmental whistleblower status, so termination is likely to trigger a lawsuit/whistleblower retaliation claim.
\item The meeting was specifically convened because outside counsel was concerned about the retaliation claim, and even some of the `ops' sections that are supposed to be nonprivileged reflect counsel's advice, so they cannot be cleanly treated as nonlegal business discussions.
\item While partially privileged might be defensible, my preferred answer would be fully privileged, which the rubric incorrectly fails.
\end{itemize}

\labitem{3}{Compare Document Production Against Discovery Requests — Discovery Gap Analysis Memorandum}{C-029}
\begin{labquote}{Harvey LAB, criterion C-029, \href{https://github.com/harveyai/harvey-labs/blob/1dd81403b2fbb60596f7aea3fcecafad7bf73143/tasks/litigation-dispute-resolution/compare-document-production-against-discovery-requests/task.json}{task file at commit \texttt{1dd81403}}}
\textbf{ISSUE\_011: Notes only 5 documents actually relate to counterclaim damages.} PASS if the memorandum notes that only 5 of the 14 documents coded to RFP 27 (MLC-004056 through MLC-004102) actually relate to Pacific Corridor's volume shortfall and the counterclaim damages calculation, meaning the counterclaim damages support is thin. FAIL if this is not noted.
\end{labquote}

\noindent\begin{tabular}{@{}>{\raggedright\arraybackslash}p{.25\linewidth}>{\raggedright\arraybackslash}p{.71\linewidth}@{}}
\toprule
\multicolumn{2}{@{}l@{}}{\textbf{My assessment}}\\
Criterion & Defective --- unambiguous objective error, rubric is incorrect\\
Task environment & Defective: the production index codes 15 rows to RFP 27, but its summary tab gives a count of 14 for a range that contains 13.\\
AI auditors' reasoning & Correct\\
\midrule
\multicolumn{2}{@{}l@{}}{\textbf{LAB's own grading of two model runs} (Harvey's judges, not my assessment)}\\
GPT-6 Luna & Sonnet 4.6 judge: fail; GPT-5.5 judge: pass\\
Claude Opus 5.5 & Sonnet 4.6 judge: fail; GPT-5.5 judge: fail\\
\bottomrule
\end{tabular}

\noindent\textbf{My analysis.} This task requires models to make the point that “only” five documents in the bundle “actually relate to Pacific Corridor's volume shortfall and the counterclaim damages calculation, meaning the counterclaim damages support is thin.” This rubric is flawed on multiple levels.

\begin{itemize}
\item First, it overlooks MLC-003781, a sixth document that plainly relates to the volume shortfall, so the count is objectively wrong (and the agents pointed to several other examples, at least some of which also seem to fit within the stated criteria). There is a document in the bundle that claims that only five documents concern the volume shortfall, but that assertion is undercut by other evidence supplied in the environment.
\item Second, this rubric is someone’s argumentative characterization of the documents and probably would not hold up under adversarial testing. For example, there is case law that the term “relate to” is fairly broad and flexible, so quite plausibly a number of other documents would apply.
\item Third, the premise underlying this rubric–that “only” five documents is such “thin” evidence that it must be called out–is wrong. Even one document could constitute conclusive proof of damages (the general rule is that damages only need to be reasonably estimated; courts do not require great precision); to my mind, this is actually quite a lot of evidence at the fact stage compared to what one would normally see.
\end{itemize}

\labitem{4}{Draft Case Assessment Memorandum — Litigation Risk Analysis for Distribution Agreement Dispute}{C-025}
\begin{labquote}{Harvey LAB, criterion C-025, \href{https://github.com/harveyai/harvey-labs/blob/1dd81403b2fbb60596f7aea3fcecafad7bf73143/tasks/litigation-dispute-resolution/draft-case-assessment-memorandum/task.json}{task file at commit \texttt{1dd81403}}}
\textbf{ISSUE\_007: Recommends motion to dismiss on unjust enrichment.} PASS if the memo recommends filing a motion to dismiss (or motion for judgment on the pleadings, or similar dispositive motion) targeting the unjust enrichment claim (Count IV) on the ground that an express contract governs the parties' relationship. FAIL if no such motion recommendation is made regarding the unjust enrichment claim.
\end{labquote}

\noindent\begin{tabular}{@{}>{\raggedright\arraybackslash}p{.25\linewidth}>{\raggedright\arraybackslash}p{.71\linewidth}@{}}
\toprule
\multicolumn{2}{@{}l@{}}{\textbf{My assessment}}\\
Criterion & Defective --- low-quality professional judgment, rubric is incorrect\\
Task environment & Correct\\
AI auditors' reasoning & Correct\\
\midrule
\multicolumn{2}{@{}l@{}}{\textbf{LAB's own grading of two model runs} (Harvey's judges, not my assessment)}\\
GPT-6 Luna & Sonnet 4.6 judge: fail; GPT-5.5 judge: fail\\
Claude Opus 5.5 & Sonnet 4.6 judge: pass; GPT-5.5 judge: pass\\
\bottomrule
\end{tabular}

\noindent\textbf{My analysis.} A different rubric on this task concludes that plaintiff has a strong breach of contract claim. (See id: C-006.) The prompt instructs models to review case documents and “prepare a comprehensive case assessment memo,” and models are dinged on this rubric if they do not “recommend[] filing a motion to dismiss” on the unjust enrichment claim, which is expressly pled in the alternative to the (strong) breach of contract claim.

\begin{itemize}
\item The rubric is wrong because a careful litigator could reasonably conclude that a motion to dismiss that claim is not a good idea: Count IV is pleaded in the alternative (Complaint ¶127) and seeks restitution of “no less than \$22,521,312.50” (¶135), which is exactly the Count I contract damages (\$17,721,312.50, ¶94) plus the Count III trade-secret damages (\$4,800,000, ¶120). Counts I and III would survive the motion and double recovery is barred, so dismissing Count IV would not reduce the defendant’s net exposure; and the factual underpinnings of both claims are the same so a successful motion would not reduce the scope of discovery. Filing a dispositive motion is costly and, on these facts, there seems to be no benefit to doing so. So the required recommendation would fail many thoughtful analyses that could discuss the pros and cons and recommend against filing such a motion.
\end{itemize}

\labitem{5}{Draft Conflict Check Memorandum for Litigation Engagement Clearance}{C-030}
\begin{labquote}{Harvey LAB, criterion C-030, \href{https://github.com/harveyai/harvey-labs/blob/1dd81403b2fbb60596f7aea3fcecafad7bf73143/tasks/litigation-dispute-resolution/draft-conflict-check-memorandum/task.json}{task file at commit \texttt{1dd81403}}}
\textbf{ISSUE\_005: Recommends screening or removal of Caleb Strand.} PASS if the memo recommends that Caleb Strand be screened from the engagement or removed from the proposed staffing. FAIL if no such recommendation is made.
\end{labquote}

\noindent\begin{tabular}{@{}>{\raggedright\arraybackslash}p{.25\linewidth}>{\raggedright\arraybackslash}p{.71\linewidth}@{}}
\toprule
\multicolumn{2}{@{}l@{}}{\textbf{My assessment}}\\
Criterion & Defective --- low-quality professional judgment, rubric is incorrect\\
Task environment & Correct\\
AI auditors' reasoning & Correct\\
\midrule
\multicolumn{2}{@{}l@{}}{\textbf{LAB's own grading of two model runs} (Harvey's judges, not my assessment)}\\
GPT-6 Luna & Sonnet 4.6 judge: pass; GPT-5.5 judge: pass\\
Claude Opus 5.5 & Sonnet 4.6 judge: pass; GPT-5.5 judge: pass\\
\bottomrule
\end{tabular}

\noindent\textbf{My analysis.} The rubric requires the model to recommend that an associate who shares an apartment with his sister (an in-house IP paralegal at TriPoint, the adverse party) be screened from the matter. Screening is generally used in certain specific situations where required by the professional conduct rules, but no rule requires screening in this scenario.

\begin{itemize}
\item I agree that a prudent memo would flag the fact that a junior team member is sharing an apartment with an employee of the adverse party, but that is covered by other rubrics (C-026 through C-029). This rubric specifically fails the work product (and results in a 0 for all work on this task) even if the model recommended, say, obtaining the client’s informed consent, which would be perfectly permissible and appropriate here.
\item Other rubrics in this task resolve similar kinds of conflicts with informed consent (C-024 accepts ``clearance with conditions (such as obtaining informed consent)'' for Lisa Chow's spousal conflict, and C-016 recommends informed consent for the Ridgeline conflict). Accordingly, informed consent plus a home confidentiality protocol could suffice as an alternative to screening.
\end{itemize}

\labitem{6}{Draft Conflict Check Memorandum for Litigation Engagement Clearance}{C-049}
\begin{labquote}{Harvey LAB, criterion C-049, \href{https://github.com/harveyai/harvey-labs/blob/1dd81403b2fbb60596f7aea3fcecafad7bf73143/tasks/litigation-dispute-resolution/draft-conflict-check-memorandum/task.json}{task file at commit \texttt{1dd81403}}}
\textbf{ISSUE\_011: Recommends integrating annual disclosure questionnaires with ConflictTracker.} PASS if the memo recommends that annual conflict-disclosure questionnaires be cross-referenced with or integrated into the ConflictTracker system. FAIL if no recommendation about integrating disclosure questionnaires with the conflicts database is made.
\end{labquote}

\noindent\begin{tabular}{@{}>{\raggedright\arraybackslash}p{.25\linewidth}>{\raggedright\arraybackslash}p{.71\linewidth}@{}}
\toprule
\multicolumn{2}{@{}l@{}}{\textbf{My assessment}}\\
Criterion & Defective --- requires specific outputs not expressly or impliedly requested in the prompt\\
Task environment & Correct\\
AI auditors' reasoning & Correct\\
\midrule
\multicolumn{2}{@{}l@{}}{\textbf{LAB's own grading of two model runs} (Harvey's judges, not my assessment)}\\
GPT-6 Luna & Sonnet 4.6 judge: fail; GPT-5.5 judge: pass\\
Claude Opus 5.5 & Sonnet 4.6 judge: pass; GPT-5.5 judge: pass\\
\bottomrule
\end{tabular}

\noindent\textbf{My analysis.} The assignment in the prompt is to draft a “comprehensive conflict check memorandum for the proposed engagement.” The rubric requires that the memo contain suggested firm-wide process improvements (memo must “recommend[] that annual conflict-disclosure questionnaires be cross-referenced with or integrated into the ConflictTracker system”). That’s simply not the assignment, and there’s many reasons why it’s undesirable to jam broad process suggestions into a conflicts memo for one specific case. This rubric actually failed work product from Luna that did everything a clearance memo needs on this point. The two judges split on that same Luna memo (Sonnet 4.6 failed it; GPT-5.5 passed it), which shows the criterion does not grade reliably.

\labitem{7}{Draft Counterclaim Against Plaintiff for Breach of Joint Development Agreement}{C-045}
\begin{labquote}{Harvey LAB, criterion C-045, \href{https://github.com/harveyai/harvey-labs/blob/1dd81403b2fbb60596f7aea3fcecafad7bf73143/tasks/litigation-dispute-resolution/draft-counterclaim-against-plaintiff-for-breach-of-joint-development-agreement/task.json}{task file at commit \texttt{1dd81403}}}
\textbf{Count: Patent Infringement (ISSUE\_005).} PASS if the counterclaim includes a separate count for patent infringement under 35 U.S.C. § 271, alleging that the LumiSense 400 infringes U.S. Patent No. 11,234,567 covering the SensorCore\texttrademark{} architecture. FAIL if no patent infringement counterclaim is asserted.
\end{labquote}

\noindent\begin{tabular}{@{}>{\raggedright\arraybackslash}p{.25\linewidth}>{\raggedright\arraybackslash}p{.71\linewidth}@{}}
\toprule
\multicolumn{2}{@{}l@{}}{\textbf{My assessment}}\\
Criterion & Arguable\\
Task environment & Defective: the record lacks the patent the counterclaim must plead, and the two expert reports give different titles for it.\\
AI auditors' reasoning & Correct\\
\midrule
\multicolumn{2}{@{}l@{}}{\textbf{LAB's own grading of two model runs} (Harvey's judges, not my assessment)}\\
GPT-6 Luna & Sonnet 4.6 judge: pass; GPT-5.5 judge: pass\\
Claude Opus 5.5 & Sonnet 4.6 judge: pass; GPT-5.5 judge: pass\\
\bottomrule
\end{tabular}

\noindent\textbf{My analysis.} The task is to draft a counterclaim based on a handful of documents (including an expert report, some emails, and a counterclaim strategy memo that another lawyer prepared). The rubric dings the AI if the counterclaims do not include “a separate count for patent infringement under 35 U.S.C. § 271, alleging that the LumiSense 400 infringes U.S. Patent No. 11,234,567.”

\begin{itemize}
\item The environment does not contain the allegedly infringed patent, which one would need to review in order to have a sufficient basis for drafting that counterclaim, so this environment appears to be incomplete. That said, a defensible approach might be for the agent to draft the counterclaim but flag in brackets that it requires verification.
\end{itemize}

\labitem{8}{Draft Counterclaim Against Plaintiff for Breach of Joint Development Agreement}{C-046}
\begin{labquote}{Harvey LAB, criterion C-046, \href{https://github.com/harveyai/harvey-labs/blob/1dd81403b2fbb60596f7aea3fcecafad7bf73143/tasks/litigation-dispute-resolution/draft-counterclaim-against-plaintiff-for-breach-of-joint-development-agreement/task.json}{task file at commit \texttt{1dd81403}}}
\textbf{Patent Infringement Count: Patent Number Cited.} PASS if the patent infringement count specifically identifies U.S. Patent No. 11,234,567. FAIL if the patent number is omitted or incorrect in the patent infringement count.
\end{labquote}

\noindent\begin{tabular}{@{}>{\raggedright\arraybackslash}p{.25\linewidth}>{\raggedright\arraybackslash}p{.71\linewidth}@{}}
\toprule
\multicolumn{2}{@{}l@{}}{\textbf{My assessment}}\\
Criterion & Arguable\\
Task environment & Defective: same task environment as item 7.\\
AI auditors' reasoning & Partly correct\\
\midrule
\multicolumn{2}{@{}l@{}}{\textbf{LAB's own grading of two model runs} (Harvey's judges, not my assessment)}\\
GPT-6 Luna & Sonnet 4.6 judge: pass; GPT-5.5 judge: pass\\
Claude Opus 5.5 & Sonnet 4.6 judge: pass; GPT-5.5 judge: pass\\
\bottomrule
\end{tabular}

\noindent\textbf{My analysis.} Same issue as prior rubric. This one is essentially derivative because it depends on pleading the patent count that was required in the prior rubric.

\labitem{9}{Draft Case Assessment Memorandum for Defective Industrial Equipment Product Liability Claim}{C-003}
\begin{labquote}{Harvey LAB, criterion C-003, \href{https://github.com/harveyai/harvey-labs/blob/1dd81403b2fbb60596f7aea3fcecafad7bf73143/tasks/litigation-dispute-resolution/draft-defective-industrial-equipment-product-liability/task.json}{task file at commit \texttt{1dd81403}}}
\textbf{Executive Summary mentions death of Marco Reyes.} PASS if the Executive Summary mentions the death of Marco Reyes. FAIL if the fatality is not mentioned in the Executive Summary.
\end{labquote}

\noindent\begin{tabular}{@{}>{\raggedright\arraybackslash}p{.25\linewidth}>{\raggedright\arraybackslash}p{.71\linewidth}@{}}
\toprule
\multicolumn{2}{@{}l@{}}{\textbf{My assessment}}\\
Criterion & Defective --- requires specific outputs not expressly or impliedly requested in the prompt\\
Task environment & Defective: same task environment as item 10 (Ohio state agency issuing OSHA citations).\\
AI auditors' reasoning & Correct\\
\midrule
\multicolumn{2}{@{}l@{}}{\textbf{LAB's own grading of two model runs} (Harvey's judges, not my assessment)}\\
GPT-6 Luna & Sonnet 4.6 judge: pass; GPT-5.5 judge: fail\\
Claude Opus 5.5 & Sonnet 4.6 judge: fail; GPT-5.5 judge: fail\\
\bottomrule
\end{tabular}

\noindent\textbf{My analysis.} The case involved a (fictitious) safety incident at a manufacturing plant that killed one employee and seriously injured two others. The rubric requires the fact that one employee died be discussed specifically in the “executive summary.” When we tested Opus 5.5 and GPT-6 Luna on this benchmark, both discussed the death in several places in their memo, especially the factual background section, but the executive summary consisted of an overview of the legal analysis that did not restate the fact of the death. They were failed by one or more judges.

\begin{itemize}
\item This fails perfectly competent and appropriate work product. The company CEO who retained counsel to advise on this incident is the same person who delivered news of the (fictitious) tragic death to the decedent’s widow and surviving family members, so the most senior management of the company is obviously very familiar with the incident and the agents’ approach of focusing the executive summary on the legal exposure is appropriate.
\item The two judges reached opposite results on Luna’s memo: GPT-5.5 failed it, while Sonnet 4.6 passed it only by quoting a sentence from the memo’s facts section as if it appeared in the executive summary.
\end{itemize}

\labitem{10}{Draft Case Assessment Memorandum for Defective Industrial Equipment Product Liability Claim}{C-026}
\begin{labquote}{Harvey LAB, criterion C-026, \href{https://github.com/harveyai/harvey-labs/blob/1dd81403b2fbb60596f7aea3fcecafad7bf73143/tasks/litigation-dispute-resolution/draft-defective-industrial-equipment-product-liability/task.json}{task file at commit \texttt{1dd81403}}}
\textbf{OSHA LOTO citation identified with specific violation and penalty.} PASS if the memo identifies the OSHA serious citation for LOTO deficiency under 29 CFR 1910.147(c)(4)(i) with a penalty of \$18,900. FAIL if this specific citation, regulation, or penalty amount is not referenced.
\end{labquote}

\noindent\begin{tabular}{@{}>{\raggedright\arraybackslash}p{.25\linewidth}>{\raggedright\arraybackslash}p{.71\linewidth}@{}}
\toprule
\multicolumn{2}{@{}l@{}}{\textbf{My assessment}}\\
Criterion & Not defective --- defective environment, but rubric does not fail correct work product\\
Task environment & Defective: Ohio's Division of Safety \& Hygiene cannot issue OSHA citations to this private employer. The criterion itself does not misgrade work (environment-only defect).\\
AI auditors' reasoning & Wrong\\
\midrule
\multicolumn{2}{@{}l@{}}{\textbf{LAB's own grading of two model runs} (Harvey's judges, not my assessment)}\\
GPT-6 Luna & Sonnet 4.6 judge: fail; GPT-5.5 judge: fail\\
Claude Opus 5.5 & Sonnet 4.6 judge: pass; GPT-5.5 judge: pass\\
\bottomrule
\end{tabular}

\noindent\textbf{My analysis.} Although the specific rubric being tested does not fail work product incorrectly, the environment is defective in an unintended way. The environment contains an OSHA citation from Ohio’s Division of Safety \& Hygiene. But Ohio does not have an OSHA-approved State Plan, so only federal regulators have jurisdiction to issue OSHA citations against this company. The environment does not test that so it is not being presented as a “mistake” by one’s adversary that the agents should pick up on.

\begin{itemize}
\item The rubric itself is not wrong because it simply wants a memo to address the purported citation, so an agent could write the memo while noticing that the task is obviously a fictitious evaluation with an incoherent environment.
\end{itemize}

\labitem{11}{Draft Discovery Plan Memorandum for Breach of Contract and Fraud Defense}{C-036}
\begin{labquote}{Harvey LAB, criterion C-036, \href{https://github.com/harveyai/harvey-labs/blob/1dd81403b2fbb60596f7aea3fcecafad7bf73143/tasks/litigation-dispute-resolution/draft-discovery-plan-memorandum/task.json}{task file at commit \texttt{1dd81403}}}
\textbf{Provides custodian priority rankings.} PASS if the memo provides priority rankings or tiers for custodians (e.g., high/medium/low priority, or Tier 1/Tier 2, or similar prioritization). FAIL if custodians are listed without any indication of relative priority.
\end{labquote}

\noindent\begin{tabular}{@{}>{\raggedright\arraybackslash}p{.25\linewidth}>{\raggedright\arraybackslash}p{.71\linewidth}@{}}
\toprule
\multicolumn{2}{@{}l@{}}{\textbf{My assessment}}\\
Criterion & Not defective\\
Task environment & Correct\\
AI auditors' reasoning & Wrong\\
\midrule
\multicolumn{2}{@{}l@{}}{\textbf{LAB's own grading of two model runs} (Harvey's judges, not my assessment)}\\
GPT-6 Luna & Sonnet 4.6 judge: pass; GPT-5.5 judge: pass\\
Claude Opus 5.5 & Sonnet 4.6 judge: pass; GPT-5.5 judge: pass\\
\bottomrule
\end{tabular}

\noindent\textbf{My analysis.} The AI agent is right that the rubric is really enforcing a stylistic choice, but the choice here is well-grounded in convention and an appropriate expectation for quality work product. The case file contains a court order where the judge encouraged the parties to discuss whether phased discovery would be appropriate at the Rule 26(f) conference, and there is some existing prioritization work in the environment. The rubric requires that the agent’s discovery plan include some kind of prioritization or phasing of custodians, which is a common expectation in these cases. I think it’s fair to judge anything that lacks any form of prioritization as deficient given the context.

\labitem{12}{Draft Federal Complaint for Breach of Contract and Fiduciary Duty — Placement Agent Dispute}{C-025}
\begin{labquote}{Harvey LAB, criterion C-025, \href{https://github.com/harveyai/harvey-labs/blob/1dd81403b2fbb60596f7aea3fcecafad7bf73143/tasks/litigation-dispute-resolution/draft-federal-complaint-drafting/task.json}{task file at commit \texttt{1dd81403}}}
\textbf{ISSUE\_012: Complaint requests injunctive relief.} PASS if the prayer for relief includes a request for preliminary and/or permanent injunctive relief — specifically to enforce the non-solicitation provision and prevent Graydon from continuing to divert Fund III investor prospects to competing funds. FAIL if no injunctive relief is requested.
\end{labquote}

\noindent\begin{tabular}{@{}>{\raggedright\arraybackslash}p{.25\linewidth}>{\raggedright\arraybackslash}p{.71\linewidth}@{}}
\toprule
\multicolumn{2}{@{}l@{}}{\textbf{My assessment}}\\
Criterion & Not defective\\
Task environment & Correct\\
AI auditors' reasoning & Wrong\\
\midrule
\multicolumn{2}{@{}l@{}}{\textbf{LAB's own grading of two model runs} (Harvey's judges, not my assessment)}\\
GPT-6 Luna & Sonnet 4.6 judge: pass; GPT-5.5 judge: pass\\
Claude Opus 5.5 & Sonnet 4.6 judge: pass; GPT-5.5 judge: pass\\
\bottomrule
\end{tabular}

\noindent\textbf{My analysis.} The client intake memo specifically describes an interest in seeking injunctive relief (“We may also wish to seek preliminary and permanent injunctive relief to enforce the twelve-month non-solicitation tail”), and the client’s demand letter threatens it. While it’s true that the intake memo flags the topic for discussion, the correct approach is to include the proposed prayer for relief in the draft, so the case team can see how it reads and discuss any downsides. As a practical matter, there’s unlikely to be much downside to including this prayer for relief in the complaint. (Whether it’s worth the cost of actually pursuing a preliminary injunction on this is debatable, so the rubric’s approach makes sense.)

\labitem{13}{Draft First Set of Interrogatories to Defendant Veridian Health Systems in Trade Secret Misappropriation Case}{C-008}
\begin{labquote}{Harvey LAB, criterion C-008, \href{https://github.com/harveyai/harvey-labs/blob/1dd81403b2fbb60596f7aea3fcecafad7bf73143/tasks/litigation-dispute-resolution/draft-interrogatories/task.json}{task file at commit \texttt{1dd81403}}}
\textbf{Definitions section defines 'Document' or 'Documents'.} PASS if the Definitions section includes a definition of 'Document' or 'Documents' that is broad enough to encompass electronically stored information (ESI), emails, and physical documents. FAIL if there is no definition of 'Document' or it is so narrow as to exclude electronic records.
\end{labquote}

\noindent\begin{tabular}{@{}>{\raggedright\arraybackslash}p{.25\linewidth}>{\raggedright\arraybackslash}p{.71\linewidth}@{}}
\toprule
\multicolumn{2}{@{}l@{}}{\textbf{My assessment}}\\
Criterion & Defective --- unambiguous objective error, rubric is incorrect\\
Task environment & Defective: the supplied documents mischaracterize or miscite case law (citation sweep; peripheral to this criterion).\\
AI auditors' reasoning & Partly correct\\
\midrule
\multicolumn{2}{@{}l@{}}{\textbf{LAB's own grading of two model runs} (Harvey's judges, not my assessment)}\\
GPT-6 Luna & Sonnet 4.6 judge: fail; GPT-5.5 judge: fail\\
Claude Opus 5.5 & Sonnet 4.6 judge: pass; GPT-5.5 judge: pass\\
\bottomrule
\end{tabular}

\noindent\textbf{My analysis.} The rubric fails any draft that does not supply a definition of “document” encompassing electronically stored information (ESI). There are two problems:

\begin{itemize}
\item First, these are interrogatories, not document requests. Well-drafted interrogatories do not necessarily need to use the word “document.” The rubric does not really make sense as a requirement for this kind of work product.
\item Second, even if we were dealing with document requests here (which we are not), the local rules of many courts, including the Western District of Texas (where this fictitious case is being litigated), define the term “document” for all discovery requests served by litigants in those cases. There is thus no requirement to define the term “documents,” and in fact, supplying a definition could just trigger objections from the other side that the local rule controls. See Local Rule CV-26, available at \url{https://www.txwd.uscourts.gov/court-information/lcr-civil-rules/} (“The term ‘document’ means any document or electronically stored information as described in Federal Rule of Civil Procedure 34(a). A draft of a nonidentical copy is a separate document within the meaning of this term.”).
\item The rubric does demonstrably fail rules-compliant work product. For example, Luna was failed for not including a definition of the term “Document” but that’s actually what the local rules encourage in this Court.
\end{itemize}

\labitem{14}{Draft First Set of Interrogatories to Defendant Veridian Health Systems in Trade Secret Misappropriation Case}{C-026}
\begin{labquote}{Harvey LAB, criterion C-026, \href{https://github.com/harveyai/harvey-labs/blob/1dd81403b2fbb60596f7aea3fcecafad7bf73143/tasks/litigation-dispute-resolution/draft-interrogatories/task.json}{task file at commit \texttt{1dd81403}}}
\textbf{Interrogatory re: USB drive or CMT files on Veridian systems.} PASS if at least one interrogatory asks whether Oshiro's USB drive (referencing serial number CX-8827491 or otherwise describing the external storage device) was ever connected to, accessed from, or present on any Veridian system, device, or network, or asks whether any CMT files or materials were ever possessed, accessed, viewed, or copied by Veridian or its employees. FAIL if no interrogatory addresses the USB drive or CMT files on Veridian systems.
\end{labquote}

\noindent\begin{tabular}{@{}>{\raggedright\arraybackslash}p{.25\linewidth}>{\raggedright\arraybackslash}p{.71\linewidth}@{}}
\toprule
\multicolumn{2}{@{}l@{}}{\textbf{My assessment}}\\
Criterion & Not defective\\
Task environment & Defective: same task environment as item 13 (environment-only defect).\\
AI auditors' reasoning & Wrong\\
\midrule
\multicolumn{2}{@{}l@{}}{\textbf{LAB's own grading of two model runs} (Harvey's judges, not my assessment)}\\
GPT-6 Luna & Sonnet 4.6 judge: pass; GPT-5.5 judge: pass\\
Claude Opus 5.5 & Sonnet 4.6 judge: pass; GPT-5.5 judge: pass\\
\bottomrule
\end{tabular}

\noindent\textbf{My analysis.} This one does not really belong on the list, as GPT-6 Sol was complaining about a different rubric (C-048) which it argues is inconsistent with this one. It does not identify any issue in this rubric, standing alone, so this counts as an instance where Sol added something to the “flagged” list that does not belong.

\labitem{15}{Draft Discovery Responses and Objections to RFAs and RFPs in Breach of Supply Agreement Litigation}{C-028}
\begin{labquote}{Harvey LAB, criterion C-028, \href{https://github.com/harveyai/harvey-labs/blob/1dd81403b2fbb60596f7aea3fcecafad7bf73143/tasks/litigation-dispute-resolution/draft-litigation-discovery-responses/task.json}{task file at commit \texttt{1dd81403}}}
\textbf{ISSUE\_011: Identifies RFA No. 20 as calling for a legal conclusion.} PASS if the plan identifies that RFA No. 20 (asking Prismavale to admit the products were 'not merchantable as defined by the Uniform Commercial Code') calls for a legal conclusion regarding UCC § 2-314 merchantability and recommends objecting on that basis. FAIL if the plan does not identify the legal conclusion problem with RFA No. 20.
\end{labquote}

\noindent\begin{tabular}{@{}>{\raggedright\arraybackslash}p{.25\linewidth}>{\raggedright\arraybackslash}p{.71\linewidth}@{}}
\toprule
\multicolumn{2}{@{}l@{}}{\textbf{My assessment}}\\
Criterion & Defective --- low-quality professional judgment, rubric is incorrect\\
Task environment & Defective: the supplied Answer cites complaint paragraph ranges that do not match the complaint (e.g. Count III as ¶¶ 88–94; the complaint's Count III is ¶¶ 86–96).\\
AI auditors' reasoning & Correct\\
\midrule
\multicolumn{2}{@{}l@{}}{\textbf{LAB's own grading of two model runs} (Harvey's judges, not my assessment)}\\
GPT-6 Luna & Sonnet 4.6 judge: fail; GPT-5.5 judge: fail\\
Claude Opus 5.5 & Sonnet 4.6 judge: pass; GPT-5.5 judge: pass\\
\bottomrule
\end{tabular}

\noindent\textbf{My analysis.} Plaintiff served a request for admission on the defendant that says, “Admit that the SB-102 and PS-302 products shipped by Prismavale to Greenleaf were not merchantable as defined by the Uniform Commercial Code.” The prompt asks the agent to “prepare a discovery strategy memo and formal RFA/RFP responses.” The rubric requires “the plan” (i.e., the strategy memo) to recommend objecting to the request on the grounds that it calls for a legal conclusion. The rubric is wrong for two main reasons:

\begin{itemize}
\item First, the modern view is requests for admissions can call for application of law to facts. While it is within acceptable professional norms for a lawyer to object on the ground that such a request calls for a legal conclusion, and then proceed to answer, the better reading of the law is that this request is not objectionable such that the party should simply answer. On this legal conclusion in particular, there’s really no reason why one would prefer to object rather than simply deny the requested admission. So this rubric fails competent work product (and probably the better answer than the one the rubric calls for).
\item Second, even assuming one should include the objection at all, the task fails the model if the issue is not expressly addressed in the strategy memo. Luna made the “legal conclusion” objection in the draft responses and then proceeded to a substantive denial (an acceptable approach), but did not discuss the “legal conclusion” objection in the memo. That was judged a failure. That’s really requiring duplicative responses on a very minor point that probably shouldn’t be in the drafts at all.
\end{itemize}

\labitem{16}{Draft Motion in Limine to Exclude Expert Testimony and Prejudicial Evidence in Commercial Breach of Contract Case}{C-009}
\begin{labquote}{Harvey LAB, criterion C-009, \href{https://github.com/harveyai/harvey-labs/blob/1dd81403b2fbb60596f7aea3fcecafad7bf73143/tasks/litigation-dispute-resolution/draft-motion-in-limine-to-exclude-expert-testimony-and-prejudicial-evidence/task.json}{task file at commit \texttt{1dd81403}}}
\textbf{ISSUE\_004: Identifies the \$11,000 arithmetic discrepancy.} PASS if the motion identifies that the sum of Marchetti's five itemized damages categories totals \$14,211,000, but her report states total damages as \$14,200,000, creating an \$11,000 discrepancy. FAIL if this arithmetic error is not identified.
\end{labquote}

\noindent\begin{tabular}{@{}>{\raggedright\arraybackslash}p{.25\linewidth}>{\raggedright\arraybackslash}p{.71\linewidth}@{}}
\toprule
\multicolumn{2}{@{}l@{}}{\textbf{My assessment}}\\
Criterion & Defective --- low-quality professional judgment, rubric is incorrect (a judgment call, not an objective error); defective environment (directs advocacy outside professional norms); companion rubric C-010 also incorrect (not sampled)\\
Task environment & Defective: the task directs advocacy outside professional norms (an \$11,000 rounding point offered as a \emph{Daubert} ground) and labels a rounding difference an ``arithmetic error.''\\
AI auditors' reasoning & Correct\\
\midrule
\multicolumn{2}{@{}l@{}}{\textbf{LAB's own grading of two model runs} (Harvey's judges, not my assessment)}\\
GPT-6 Luna & Sonnet 4.6 judge: pass; GPT-5.5 judge: pass\\
Claude Opus 5.5 & Sonnet 4.6 judge: pass; GPT-5.5 judge: pass\\
\bottomrule
\end{tabular}

\noindent\textbf{My analysis.} The damages expert includes a summary table in her report listing various items of damages, which collectively add up to \$14,211,000; the Total Damages line lists the total as \$14,200,000. Opposing counsel questioned the expert about it at her deposition, and the expert calls it “a rounding discrepancy.” An internal strategy memo highlights the issue as a ground for seeking to exclude the expert’s testimony under \emph{Daubert}, arguing “the arithmetic error supports a cumulative argument that the report was prepared without the rigor required of expert testimony in federal court.”

\begin{itemize}
\item This task is poorly specified, and the recommended strategy of raising this with the judge is a poor one:
\item First, the task environment repeatedly states that this is an “arithmetic error.” That is false. The expert has rounded the total in the table to the nearest \$100,000. While that is somewhat unorthodox in an expert report (particularly since it is written out as \$14,200,000 not \$14.2 million), numbers often are rounded in court filings so the judge is unlikely to agree that this is an error.
\item Second, even if there were an error (which there is not), that is generally not a basis for excluding the expert’s testimony under \emph{Daubert}. At most one could move to exclude the erroneous version and ask for more accurate testimony, but that does nothing to help the hypothetical client here; it just raises their damages exposure by the \$11,000 that the expert was willing to round off.
\item Third, the proposed framing comes close to a breach of decorum and carries a real risk of drawing a rebuke from the judge. The strategy memo wants the motion to argue the expert’s 30 prior engagements were “in state courts and arbitrations that do not impose the same methodological gatekeeping requirements” as federal courts. While some state courts use the \emph{Frye} standard instead of \emph{Daubert}, many federal judges would take offense at an argument that state courts do not observe adequate methodological rigor. And the reality is that many busy judges would be annoyed at spending time on an argument that literally amounts to a rounding error. It's so trivial that raising it as a substantive point risks making the judge think you had no better argument to put forward, so your motion must be weak.
\item This rubric is scored as Defective, as a professional-judgment call rather than an objective error. Our view is that the better approach is to drop this issue entirely, and the sampled rubric (C-009) fails a motion that does so, even though a motion that merely mentions the discrepancy in a footnote would pass. The companion rubric, C-010, is in my view objectively wrong because it requires that this issue be presented as a basis for excluding the expert's testimony under \emph{Daubert}. Luna was failed incorrectly on C-010 (by the Sonnet 4.6 judge) for downplaying this issue; in fact, Luna was right to do so, as it should not be raised at all. But since that's not the sampled criterion, we do not mark it here.
\end{itemize}

\labitem{17}{Draft Rule 12(b)(6) Motion to Dismiss Brief — Commercial Software Licensing Dispute}{C-037}
\begin{labquote}{Harvey LAB, criterion C-037, \href{https://github.com/harveyai/harvey-labs/blob/1dd81403b2fbb60596f7aea3fcecafad7bf73143/tasks/litigation-dispute-resolution/draft-motion-to-dismiss-brief/task.json}{task file at commit \texttt{1dd81403}}}
\textbf{ISSUE\_008 — Cites Delaware parol evidence authority.} PASS if the motion brief cites relevant Delaware authority on the parol evidence rule and integration clauses, such as SIGA Technologies v. PharmAthene or Eagle Industries v. DeVilbiss Health Care. FAIL if no Delaware parol evidence authority is cited.
\end{labquote}

\noindent\begin{tabular}{@{}>{\raggedright\arraybackslash}p{.25\linewidth}>{\raggedright\arraybackslash}p{.71\linewidth}@{}}
\toprule
\multicolumn{2}{@{}l@{}}{\textbf{My assessment}}\\
Criterion & Defective --- unambiguous objective error (including hallucinations), rubric is incorrect\\
Task environment & Defective: the supplied research memo cites fabricated or mischaracterized authority (citation sweep: 28 of 53 citations in this task have problems).\\
AI auditors' reasoning & Correct\\
\midrule
\multicolumn{2}{@{}l@{}}{\textbf{LAB's own grading of two model runs} (Harvey's judges, not my assessment)}\\
GPT-6 Luna & Sonnet 4.6 judge: fail; GPT-5.5 judge: fail\\
Claude Opus 5.5 & Sonnet 4.6 judge: fail; GPT-5.5 judge: fail\\
\bottomrule
\end{tabular}

\noindent\textbf{My analysis.} The rubric here involves a legal error, and the task environment contains hallucinations about case law. This one is extremely defective for multiple reasons:

\begin{itemize}
\item First, the rubric is demanding case law on the wrong legal doctrine. Under Delaware law, the parol evidence rule and an integration clause do not preclude fraud claims. The issue here is whether the plaintiff can reasonably allege that they relied on their counterparty’s pre-signing representations despite their agreement in the contract that they were not relying on any representations other than those written in the contract itself.
\item Second, the cases in the environment appear to contain multiple hallucinations. The rubric refers to SIGA v. PharmAthene, but Opus noted that case does not address integration or parol evidence. The task environment contains a legal research memo that cites \emph{SIGA Technologies, Inc. v. PharmAthene, Inc.}, 67 A.3d 330 (Del. 2013), and claims that it states (at page 344), ``an integration clause is the clearest declaration that the parties intend the writing to be the complete and final expression of their agreement.'' The quoted language does not appear in the cited case. Page 344 of that case addresses the duty to negotiate in good faith, not integration clauses. Similarly, the memo cites \emph{Eagle Industries, Inc. v. DeVilbiss Health Care, Inc.}, 702 A.2d 1228 (Del. 1997) and says: “The Delaware Supreme Court held that ‘a party to a contract cannot promise, in a clear integration clause of a negotiated agreement, that it is not relying on promises or representations made by the other party outside of the agreement, and then assert a claim for fraud based on those very representations.’ \emph{Id.} at 1232.” The quoted language does not appear in the cited case.
\item Third, the rubric failed responses that stated the correct law on this point. In our tests, Opus reviewed SIGA and Eagle, correctly realized they were inapposite and declined to cite them, cited more on-point and accurate authority, and was then failed for doing a better job than the rubric's author did at this task.
\end{itemize}

\labitem{18}{Draft Rule 12(b)(6) Motion to Dismiss Brief — Commercial Software Licensing Dispute}{C-045}
\begin{labquote}{Harvey LAB, criterion C-045, \href{https://github.com/harveyai/harvey-labs/blob/1dd81403b2fbb60596f7aea3fcecafad7bf73143/tasks/litigation-dispute-resolution/draft-motion-to-dismiss-brief/task.json}{task file at commit \texttt{1dd81403}}}
\textbf{ISSUE\_010 — References Arcadia's own delays contributing to Go-Live delay.} PASS if the motion brief identifies Arcadia's own contributing failures, including the 45-day Project Manager gap (Kevin Liu's departure) and/or the 75-day late API specifications, as causes of the Go-Live delay. FAIL if Arcadia's own contributory delays are not mentioned.
\end{labquote}

\noindent\begin{tabular}{@{}>{\raggedright\arraybackslash}p{.25\linewidth}>{\raggedright\arraybackslash}p{.71\linewidth}@{}}
\toprule
\multicolumn{2}{@{}l@{}}{\textbf{My assessment}}\\
Criterion & Defective --- unambiguous objective error, rubric is incorrect\\
Task environment & Defective: same task environment as item 17.\\
AI auditors' reasoning & Correct\\
\midrule
\multicolumn{2}{@{}l@{}}{\textbf{LAB's own grading of two model runs} (Harvey's judges, not my assessment)}\\
GPT-6 Luna & Sonnet 4.6 judge: fail; GPT-5.5 judge: fail\\
Claude Opus 5.5 & Sonnet 4.6 judge: fail; GPT-5.5 judge: fail\\
\bottomrule
\end{tabular}

\noindent\textbf{My analysis.} The rubric makes a rookie mistake, demanding that the agent argue a factual dispute on a motion to dismiss, where the court is required to assume the plaintiff’s allegations are true and cannot consider factual dispute. The 45-day and 75-day delays complained of come from the defendant’s internal report, which the plaintiff’s complaint disputes. Moreover, courts generally cannot resolve messy fact disputes over “causation” on a motion to dismiss.

\begin{itemize}
\item The criterion failed agents that correctly recognized that this argument was outside the scope of what can be argued on a motion to dismiss and thus decided to exclude it.
\end{itemize}

\labitem{19}{Draft Verified Responses and Objections to Plaintiff's First Set of Interrogatories in Commercial Breach of Contract and Fraud Action}{C-028}
\begin{labquote}{Harvey LAB, criterion C-028, \href{https://github.com/harveyai/harvey-labs/blob/1dd81403b2fbb60596f7aea3fcecafad7bf73143/tasks/litigation-dispute-resolution/draft-responses-to-interrogatories/task.json}{task file at commit \texttt{1dd81403}}}
\textbf{Correct financial figures for Tri-Basin purchase history.} PASS if any purchase history figures cited in the responses are consistent with the canonical data: 2015: \$9.2M; 2016: \$10.1M; 2017: \$10.8M; 2018: \$12.4M; 2019: \$13.6M; 2020: \$9.7M; 2021: \$15.3M; 2022: \$16.8M; 2023 (Jan–Aug): \$11.9M. Minor rounding differences are acceptable. Also PASS if the response invokes Rule 33(d) rather than citing specific figures. FAIL if specific purchase figures are cited but are materially incorrect.
\end{labquote}

\noindent\begin{tabular}{@{}>{\raggedright\arraybackslash}p{.25\linewidth}>{\raggedright\arraybackslash}p{.71\linewidth}@{}}
\toprule
\multicolumn{2}{@{}l@{}}{\textbf{My assessment}}\\
Criterion & Defective --- unambiguous objective error, deficiency in task environment; rubric is incorrect\\
Task environment & Defective: the purchase-history spreadsheet is internally inconsistent (three 2020 series; line items that do not sum to subtotals), and the record has no contract-year data.\\
AI auditors' reasoning & Partly correct\\
\midrule
\multicolumn{2}{@{}l@{}}{\textbf{LAB's own grading of two model runs} (Harvey's judges, not my assessment)}\\
GPT-6 Luna & Sonnet 4.6 judge: pass; GPT-5.5 judge: pass\\
Claude Opus 5.5 & Sonnet 4.6 judge: pass; GPT-5.5 judge: pass\\
\bottomrule
\end{tabular}

\noindent\textbf{My analysis.} There are two issues:

\begin{itemize}
\item First, Interrogatory No. 4(b) asks for Tri-Basin’s actual purchases for the 2020 contract year (the contract year runs from January 15 through January 14 of the following year; Interrogatory No. 14 separately asks for calendar-year totals), but the canonical data referred to in the rubric is calendar year data, so the rubric's numbers are not correct for that request. The correct value cannot be discerned from the environment: the monthly data cannot be split at January 15, and nothing in the record supports assuming purchases are spread evenly within January (seasonal variation would make that unlikely). The rubric therefore grades wrong in the opposite direction from most of the other items: rather than failing good work product, it passes responses that supply the calendar-year figure for a contract-year request without acknowledging that the record does not contain the requested data. That is approving fundamentally deficient work product.
\item Second, the data that is provided has some internal inconsistencies. For example, an “annual summary” tab lists 2020 purchases at \$9.7 million, but the monthly detail lists 2020 three times, with totals of \$8,700,000, \$8,549,400 and \$9,700,000. The units sold do not match up between those two sheets for 2021, 2022, 2023, and series 3 of 2020. One of the monthly rows does not add up (gross \$1,210,000 + credits $-$\$17,400 = \$1,192,600, but the net shown is \$1,192,500). The annual summary minimums do not add up. In the Product Line Breakdown, line-item net revenue never sums to the subtotal. So a response that relied on the line items would get graded as a failure.
\item The environment contains law firm discovery guidelines that recommend that interrogatory responses match the client's Answer, but the Answer here contains only approximate figures, which is consistent with the Answer using calendar year data as a rough approximation instead of providing the contract year data requested.
\end{itemize}

\labitem{20}{Extract Key Obligations from Litigation Hold and Document Preservation Notice — Obligation Summary Memorandum}{C-002}
\begin{labquote}{Harvey LAB, criterion C-002, \href{https://github.com/harveyai/harvey-labs/blob/1dd81403b2fbb60596f7aea3fcecafad7bf73143/tasks/litigation-dispute-resolution/extract-key-obligations-from-litigation-hold-and-document-preservation-notice/task.json}{task file at commit \texttt{1dd81403}}}
\textbf{ISSUE\_001: Flags financial records date discrepancy as requiring DOJ clarification.} PASS if the memo recommends seeking clarification from the DOJ regarding the conflicting date ranges (Jan 1, 2017 vs. Jan 1, 2019) for financial records. FAIL if no recommendation to seek clarification is made.
\end{labquote}

\noindent\begin{tabular}{@{}>{\raggedright\arraybackslash}p{.25\linewidth}>{\raggedright\arraybackslash}p{.71\linewidth}@{}}
\toprule
\multicolumn{2}{@{}l@{}}{\textbf{My assessment}}\\
Criterion & Defective --- unambiguous objective error, rubric is incorrect\\
Task environment & Correct\\
AI auditors' reasoning & Correct\\
\midrule
\multicolumn{2}{@{}l@{}}{\textbf{LAB's own grading of two model runs} (Harvey's judges, not my assessment)}\\
GPT-6 Luna & Sonnet 4.6 judge: fail; GPT-5.5 judge: fail\\
Claude Opus 5.5 & Sonnet 4.6 judge: pass; GPT-5.5 judge: fail\\
\bottomrule
\end{tabular}

\noindent\textbf{My analysis.} The task involves a document preservation notice from federal criminal authorities that defines the relevant period (for which the recipient must preserve documents) as January 1, 2019, through March 3, 2025 “unless otherwise specified herein.” For one specific category of documents, the notice states that, “Notwithstanding the Relevant Period defined in Paragraph 12 of this Notice, Ridgeline shall preserve [certain kinds of] records . . . from January 1, 2017 to the present date.”

\begin{itemize}
\item The rubric suggests that the agent should seek clarification of the “conflicting” date ranges but there is no conflict. The notice quite clearly imposed a longer period for one specific category of documents, while all the other categories are subject to a shorter period. This is a fairly common request.
\item All the models were failed on this criterion, but the criterion makes no sense.
\end{itemize}

\labitem{21}{Identify Excessive or Duplicative Research Charges in Litigation Invoice}{C-013}
\begin{labquote}{Harvey LAB, criterion C-013, \href{https://github.com/harveyai/harvey-labs/blob/1dd81403b2fbb60596f7aea3fcecafad7bf73143/tasks/litigation-dispute-resolution/identify-excessive-or-duplicative-research-charges-in-litigation-invoice/task.json}{task file at commit \texttt{1dd81403}}}
\textbf{Issue 3: Identifies duplicative MTCA research across Takahashi and Wendt.} PASS if the memo identifies that multiple attorneys (Kenji Takahashi and Tyler Wendt) billed for overlapping MTCA liability/contractor liability/summary judgment research during July 2024, citing Section 6.4 (Duplicative Research Prohibition). FAIL if the MTCA duplication across timekeepers is not identified.
\end{labquote}

\noindent\begin{tabular}{@{}>{\raggedright\arraybackslash}p{.25\linewidth}>{\raggedright\arraybackslash}p{.71\linewidth}@{}}
\toprule
\multicolumn{2}{@{}l@{}}{\textbf{My assessment}}\\
Criterion & Not defective\\
Task environment & Correct\\
AI auditors' reasoning & Wrong\\
\midrule
\multicolumn{2}{@{}l@{}}{\textbf{LAB's own grading of two model runs} (Harvey's judges, not my assessment)}\\
GPT-6 Luna & Sonnet 4.6 judge: pass; GPT-5.5 judge: pass\\
Claude Opus 5.5 & Sonnet 4.6 judge: pass; GPT-5.5 judge: pass\\
\bottomrule
\end{tabular}

\noindent\textbf{My analysis.} While Sol is right that the facts here are insufficient to establish that the client’s billing guidelines were violated, for purposes of the task (a client-side memo identifying changes to dispute), it seems reasonable to flag that multiple attorneys were doing research into closely related issues. The task requested “a memo identifying all non-compliant or excessive charges with recommended adjustments,” so flagging potential violations (even if they were not clearly proven from the face of the time entries) is a standard approach a client would take here.

\labitem{22}{Identify Issues in Litigation Matter Budget Proposal}{C-029}
\begin{labquote}{Harvey LAB, criterion C-029, \href{https://github.com/harveyai/harvey-labs/blob/1dd81403b2fbb60596f7aea3fcecafad7bf73143/tasks/litigation-dispute-resolution/identify-issues-in-matter-budget-proposal/task.json}{task file at commit \texttt{1dd81403}}}
\textbf{Budget inconsistency rated as Critical or High severity.} PASS if the fee total reconciliation discrepancy (\$116,050 difference between phase-based and staffing-model totals) is assigned a severity of Critical or High (or the highest or second-highest tier in whatever severity scale is used). FAIL if it is assigned a lower severity level or no severity level.
\end{labquote}

\noindent\begin{tabular}{@{}>{\raggedright\arraybackslash}p{.25\linewidth}>{\raggedright\arraybackslash}p{.71\linewidth}@{}}
\toprule
\multicolumn{2}{@{}l@{}}{\textbf{My assessment}}\\
Criterion & Arguable --- requires specific outputs not expressly or impliedly requested in the prompt\\
Task environment & Correct\\
AI auditors' reasoning & Correct\\
\midrule
\multicolumn{2}{@{}l@{}}{\textbf{LAB's own grading of two model runs} (Harvey's judges, not my assessment)}\\
GPT-6 Luna & Sonnet 4.6 judge: pass; GPT-5.5 judge: pass\\
Claude Opus 5.5 & Sonnet 4.6 judge: pass; GPT-5.5 judge: pass\\
\bottomrule
\end{tabular}

\noindent\textbf{My analysis.} The prompt requested “a categorized issue memorandum” but didn’t specify what the categories should be. The rubric fails the model if it does not break issues into severity levels or it does not assign a critical or high severity level to a specific issue. While it appears that models often do volunteer this kind of severity assessment, it’s not requested by the prompt and competent work product could decline to include it.

\labitem{23}{Review Counterparty's Proposed Jury Instructions — Issue Memorandum for Trade Secrets Trial}{C-020}
\begin{labquote}{Harvey LAB, criterion C-020, \href{https://github.com/harveyai/harvey-labs/blob/1dd81403b2fbb60596f7aea3fcecafad7bf73143/tasks/litigation-dispute-resolution/review-counterpartys-proposed-jury-instructions/task.json}{task file at commit \texttt{1dd81403}}}
\textbf{Explains distinction between tortious interference with contract vs. business relations.} PASS if the memo explains that under Georgia law, tortious interference with an existing contract does not require a separately/independently wrongful act — the knowing inducement of the breach is sufficient — and that the 'independently wrongful act' element applies to the distinct tort of tortious interference with business relations. FAIL if this legal distinction is not articulated.
\end{labquote}

\noindent\begin{tabular}{@{}>{\raggedright\arraybackslash}p{.25\linewidth}>{\raggedright\arraybackslash}p{.71\linewidth}@{}}
\toprule
\multicolumn{2}{@{}l@{}}{\textbf{My assessment}}\\
Criterion & Defective --- unambiguous objective error (including hallucinations), rubric is incorrect\\
Task environment & Defective: the fictional summary-judgment order cites a case for the opposite of its holding and a case that cannot be found (citation sweep: 10 of 21 citations in this task have problems).\\
AI auditors' reasoning & Correct\\
\midrule
\multicolumn{2}{@{}l@{}}{\textbf{LAB's own grading of two model runs} (Harvey's judges, not my assessment)}\\
GPT-6 Luna & Sonnet 4.6 judge: fail; GPT-5.5 judge: fail\\
Claude Opus 5.5 & Sonnet 4.6 judge: fail; GPT-5.5 judge: fail\\
\bottomrule
\end{tabular}

\noindent\textbf{My analysis.} Multiple issues:

\begin{itemize}
\item First, this is a clear legal error. Georgia applies the same elements to interference with contractual and business relations. The rubric says the opposite of what Georgia law is. \emph{Disaster Services, Inc. v. ERC P'ship}, 228 Ga. App. 739, 740–41 (1997) (“Tortious interference claims, whether asserting interference with contractual relations, business relations, or potential business relations, share certain common essential elements: (1) improper action or wrongful conduct by the defendant without privilege; (2) the defendant acted purposely and with malice with the intent to injure; (3) the defendant induced a breach of contractual obligations or caused a party or third parties to discontinue or fail to enter into an anticipated business relationship with the plaintiff; and (4) the defendant's tortious conduct proximately caused damage to the plaintiff.”).
\item Second, the fictional summary judgment order (from which the rubric’s legal position seems to be taken) appears to contain several hallucinations. For example, the order paraphrases the third and fourth element described in \emph{Disaster Services} approximately correctly, but the first two elements are erroneously listed as “(1) the existence of a valid contract between the plaintiff and a third party; (2) the defendant's knowledge of the contract.” The fictional order also cites \emph{Valdosta Livestock, Inc. v. Furst}, 342 Ga. App. 25, 28 (2017), which is not found in any legal databases.
\end{itemize}

\labitem{24}{Review Litigation Invoice Against Outside Counsel Billing Guidelines — Compliance Deviation Report}{C-024}
\begin{labquote}{Harvey LAB, criterion C-024, \href{https://github.com/harveyai/harvey-labs/blob/1dd81403b2fbb60596f7aea3fcecafad7bf73143/tasks/litigation-dispute-resolution/review-litigation-invoice-against-outside-counsel-billing-guidelines/task.json}{task file at commit \texttt{1dd81403}}}
\textbf{Recommends full disallowance of black car service: \$387.00.} PASS if the report recommends disallowing the full \$387.00 for black car service (or reducing it to a reasonable taxi/ride-share equivalent). FAIL if no disallowance or reduction is recommended.
\end{labquote}

\noindent\begin{tabular}{@{}>{\raggedright\arraybackslash}p{.25\linewidth}>{\raggedright\arraybackslash}p{.71\linewidth}@{}}
\toprule
\multicolumn{2}{@{}l@{}}{\textbf{My assessment}}\\
Criterion & Not defective\\
Task environment & Correct\\
AI auditors' reasoning & Wrong\\
\midrule
\multicolumn{2}{@{}l@{}}{\textbf{LAB's own grading of two model runs} (Harvey's judges, not my assessment)}\\
GPT-6 Luna & Sonnet 4.6 judge: pass; GPT-5.5 judge: pass\\
Claude Opus 5.5 & Sonnet 4.6 judge: pass; GPT-5.5 judge: pass\\
\bottomrule
\end{tabular}

\noindent\textbf{My analysis.} The billing guidelines state that “Reimbursable ground transportation is limited to ride-share \ldots{} or taxi” and “Limousine services and black car services are not reimbursable under any circumstances.” Sol’s critique of the rubric is that the policy requires disallowing black car reimbursement, not reducing the reimbursement amount to what a taxi or ride share would have cost. The rubric’s interpretation (that either disallowing the expense entirely or capping reimbursement at the amount a taxi would have cost) seems reasonable and closer to the likely real-world application of a policy like this.

\labitem{25}{Review Litigation Invoice Against Outside Counsel Billing Guidelines — Compliance Deviation Report}{C-042}
\begin{labquote}{Harvey LAB, criterion C-042, \href{https://github.com/harveyai/harvey-labs/blob/1dd81403b2fbb60596f7aea3fcecafad7bf73143/tasks/litigation-dispute-resolution/review-litigation-invoice-against-outside-counsel-billing-guidelines/task.json}{task file at commit \texttt{1dd81403}}}
\textbf{Classifies business class airfare as borderline/requires discussion.} PASS if the report classifies the business class airfare issue (Messina's \$2,847.00 flight) as a borderline or 'requires discussion' issue rather than a clear-cut mandatory reduction, given the 4-hour-20-minute flight duration and the 4-hour threshold. FAIL if it is classified as a clear mandatory reduction without acknowledging the ambiguity, OR if it is entirely dismissed as compliant.
\end{labquote}

\noindent\begin{tabular}{@{}>{\raggedright\arraybackslash}p{.25\linewidth}>{\raggedright\arraybackslash}p{.71\linewidth}@{}}
\toprule
\multicolumn{2}{@{}l@{}}{\textbf{My assessment}}\\
Criterion & Defective --- unambiguous objective error, rubric is incorrect\\
Task environment & Correct\\
AI auditors' reasoning & Correct\\
\midrule
\multicolumn{2}{@{}l@{}}{\textbf{LAB's own grading of two model runs} (Harvey's judges, not my assessment)}\\
GPT-6 Luna & Sonnet 4.6 judge: pass; GPT-5.5 judge: pass\\
Claude Opus 5.5 & Sonnet 4.6 judge: pass; GPT-5.5 judge: pass\\
\bottomrule
\end{tabular}

\noindent\textbf{My analysis.} The rubric’s analysis is wrong. The travel policy included in the environment states, “For domestic flights with a scheduled duration of four (4) hours or more, business class may be booked without prior approval.” The record does not state the flight’s duration (the invoice lists only the route and fare class). The only basis for raising this seems to be in the rubric itself, which says the flight could take 4h20m, but even under that view, there’s nothing to discuss, as the policy allows business class for that duration.
\end{document}
